\chapter{The Simple Sentence}\label{ch:simple}

% ============================================================
% STATUS: SCAFFOLD — largely new content; some material available in
% archive/2020/chapters/07-clause-structure.tex.
% Key theoretical framing: topic-comment structure is the primary
% organizational principle of the Iridian clause. The split-ergative
% alignment can be recast as: the absolutive is typically the topic;
% the ergative demotes the agent away from topic status.
% ============================================================

\section{Topic-comment structure and information structure}\label{sec:topic-comment}

% ============================================================
% The basic Iridian clause is organized around a topic-comment structure.
% The topic (what the sentence is "about") is unmarked (absolutive) and
% typically appears sentence-initially. The comment (what is said about
% the topic) follows. The alignment system is closely tied to this:
%   - Perfective/retrospective: P is the topic (ABS), A is demoted (ERG)
%   - Nominative aspects: A/S is the topic (ABS), P is non-topic (OBL)
% This means aspect choice is partly pragmatic — the speaker chooses the
% aspect that places the intended topic in the unmarked (absolutive) position.
% Antipassive demotes P and restores A to topic/ABS position.
% ============================================================

The basic Iridian clause organizes information around a topic-comment structure. The \term{topic} --- the entity the sentence is about --- is expressed in the absolutive case and typically occupies the initial position in the clause. The \term{comment} --- what is predicated of the topic --- follows. This organization is reflected in and reinforced by the split-ergative alignment: in the perfective and retrospective aspects, the patient is promoted to the absolutive (becoming the topic), while in the nominative aspects, the agent occupies the absolutive (topic) position.

Aspect choice is therefore partly a pragmatic decision: the speaker selects the aspect that places the discourse-salient argument in the absolutive position.

% [Examples to be added]

\section{Word order}\label{sec:word-order}

The canonical word order in Iridian is SOV (subject-object-verb), where the finite verb is clause-final. This order is relatively rigid for the finite verb itself, which must appear last. The ordering of nominal arguments before the verb is more flexible and is governed by information structure rather than strict grammatical requirements. Topic-initial and focus-final tendencies apply within the pre-verbal domain.

\subsection{Lexical manner and frequency adverbs in the preverbal field}
\label{sec:lexical-adverb-position}

Lexical adverbs belong to the comment domain of the clause rather than to the
topic field. In the unmarked structure, derived manner adverbs in \ird{-ě} are
best analyzed as low adjuncts adjoined to vP or to the outer VP. Because
Iridian is head-final, these adverbs surface in the preverbal field, below
topic and focus material and outside the fixed hierarchy of preverbal particles
described in Chapter~\ref{ch:function-words}. Their neutral placement is
therefore lower than that of clause-typing or stance particles but still to
the left of the finite verb.

\begin{center}
\term{[TopP Topic [IP \ldots [vP MANNER-ADV [VP Object V ]]]]}
\end{center}

\noindent The neutral surface order in an intransitive clause is accordingly
\term{Adverb $>$ Verb}:

\pex
\a
\begingl
\gla Skorně zmačime.//
\glb quickly run-\Prog{}//
\glft \trsl{I am running quickly.}//
\endgl
\xe

\noindent In a transitive clause, where the object remains in the low
complement domain, the neutral order is \term{Adverb $>$ Object $>$ Verb}:

\pex
\a
\begingl
\gla Skorně bylam piaštébime.//
\glb quickly child.\Dir{} eat-\Appl{}-\Prog{}//
\glft \trsl{I am feeding the child quickly.}//
\endgl
\xe

\noindent Frequency adverbs such as \ird{ohly} \trsl{sometimes},
\ird{roca} \trsl{often}, \ird{člíně} \trsl{infrequently}, and \ird{sabo}
\trsl{never} belong to the same broad lexical-adverb domain, though their
semantics are higher and more event-quantificational than those of pure manner
adverbs. Even so, they are not members of the rigid particle series. More
leftward placement of lexical adverbs is possible under discourse pressure, but
such ordering belongs to marked information structure and should not be taken
as the neutral pattern.

% [Examples: SOV default, pragmatic reordering, focus movement]

\section{The copula and existential constructions}\label{sec:copula}

Iridian distributes predication across three structurally and semantically distinct constructions. Where many European languages rely on a single copular verb for all such functions, Iridian treats equative identity, marked stative or inchoative predication, and existential predication separately. All overt predicative verbs are clause-final and inflect for the full aspect system; each paradigm is strongly suppletive and morphologically irregular.

Ordinary property predication belongs to the verbal system. Meanings such as
\trsl{be beautiful}, \trsl{be good}, and \trsl{be open} are normally expressed
with stative verbs, not with bare adjectival predicates. For the small closed
set of lexical property roots discussed in Chapters~\ref{ch:overview} and
\ref{ch:verbs}, the root itself is chiefly adnominal, while the corresponding
stative verb serves as the regular clause predicate. Because stative verbs
already distinguish state from entry into state aspectually, the progressive
naturally yields inchoative readings such as \trsl{becoming beautiful} or
\trsl{getting better}, while the perfective yields \trsl{became beautiful} or
\trsl{turned good}.

\subsection{Equative and class-membership clauses: the zero copula}\label{sec:zero-copula}

The default pattern for equative clauses --- those asserting identity, role, or
class membership --- is zero copula in the indicative. The predicative
complement appears in the direct case in clause-final position; no overt verb is
present. This is the stylistically neutral, unmarked construction. When such
clauses are negated, \ird{staly} appears in the clause-final position. The verbal
negative prefix \ird{zá-} is not used: it requires an overt verbal host, which
the zero-copula construction lacks.

Structurally, \ird{staly} is the phonological exponent of the NegP head
(Neg\textsuperscript{0}) in equative contexts, not a copular verb. It bears no
aspect marking, has no paradigmatic alternants, and cannot appear in any
non-negative environment. Its clause-final position is the expected one under
head-finality: as the head of NegP, it follows its complement (the nominal
predication). In verbal clauses the same NegP head is spelled out as the prefix
\ird{zá-}, which cliticizes to the verbal complex; in equative clauses, with no
verbal host available, Neg\textsuperscript{0} surfaces in its canonical
head-final position as a free morpheme. See §\,\ref{sec:clausal-negation} for
the unified treatment and §\,\ref{sec:vm-neg-exponents} (\texttt{verb-minim.tex})
for the formal analysis.

\pex
\a
\begingl
\gla Marek doktor.//
\glft \trsl{Marek is a doctor.}//
\endgl
\a
\begingl
\gla Bách němscvení kompozitor.//
\glft \trsl{Bach is a German composer.}//
\endgl

\xe

Inclusive and class-membership predications --- of the type \textit{dogs are
animals} --- follow the same zero-copula pattern and carry no inference beyond
simple predication. The predicative complement always appears in the direct case
regardless of which case would otherwise be assigned in the same position.

\pex
\a
\begingl
\gla Marek doktor staly.//
\glb Marek.\Dir{} doctor.\Dir{} \NegCop{}//
\glft \trsl{Marek is not a doctor.}//
\endgl
\a
\begingl
\gla Doktor Marek staly.//
\glb doctor.\Dir{} Marek.\Dir{} \NegCop{}//
\glft \trsl{Marek is not the doctor.}//
\endgl
\xe

\noindent The two examples differ in which DP occupies the topic (initial)
position, not in the position of \ird{staly}. Since \ird{staly} is the
clause-final NegP head, it is invariably last; the two nominal DPs may
permute before it in accordance with the general topic-comment organization
of Iridian clauses (§\,\ref{sec:topic-comment}). The semantic contrast --- not
a doctor vs.\ not the doctor --- arises from which DP is topicalized, not from
any change in the negation structure.

\ird{Staly} is restricted to the negation of equative and
class-membership clauses. It does not replace \ird{neh-}, whose overt use
remains semantically marked for becoming or entry into a state, and it does not
replace the existential negative verb \ird{zad-}, which is used for absence and
non-existence.

\subsection{Identificational equatives and participant nominals}\label{sec:ident-equative}

When the predicative complement of a zero-copula equative clause is a
\term{participant nominal} (a nominalized verb form, §\,\ref{sec:nmlz-predicative}
of Chapter~\ref{ch:verbs}), the construction is \term{identificational}: the
topic DP is asserted to be the unique referent satisfying the description encoded
in the nominal. The schematic contrast with a plain verbal assertion is:

\pex
\a
\begingl
\gla Lobek Janku piaštnik.//
\glb apple.\Dir{} Janek.\Trs{} eat-\textsc{trs}-\Pf{}//
\glft \trsl{Janek ate the apple.} \lingcontext{verbal clause: propositional assertion}//
\endgl
\a
\begingl
\gla Lobek Jancí piaštnikou.//
\glb apple.\Dir{} Janek.\Ind{} eat-\textsc{trs}-\Pf{}-\textsc{nmlz}//
\glft \trsl{The apple is what Janek ate.} \lingcontext{identificational equative: patient nominal}//
\endgl
\xe

\noindent In (b), the apple is not merely in the event but is being identified
as the eaten-thing. An agent-nominal equative identifies the topic as the agent of
the nominalized event:

\ex
\begingl
\gla Janek lubcí piaščikou.//
\glb Janek.\Dir{} apple.\Ind{} eat-\Antip{}-\Pf{}-\textsc{nmlz}//
\glft \trsl{Janek is the one who ate the apple.}//
\endgl
\xe

\noindent All internal arguments of the nominalized event appear in the indirect
case (\Ind{}) regardless of the case they would bear in the finite verbal clause:
the agent of the transitive verb (normally \Trs{}) and the patient of the
antipassive-derived agent nominal (normally suppressed in finite contexts) both
appear as \Ind{} arguments within the nominal predicate.

\paragraph{Negation of identificational equatives.}
When an identificational equative is negated with \ird{staly}, the result is
\term{identificational negation}: the topic DP is denied as the unique satisfier
of the description, with the strong implication that some other referent is the
actual satisfier. This contrasts systematically with verbal propositional negation:

\pex
\a
\begingl
\gla Lobek Janku zápiaštnik.//
\glb apple.\Dir{} Janek.\Trs{} \Neg{}-eat-\textsc{trs}-\Pf{}//
\glft \trsl{Janek did not eat the apple.} \lingcontext{propositional negation: denies that the event occurred}//
\endgl
\a
\begingl
\gla Lobek Jancí piaštnikou staly.//
\glb apple.\Dir{} Janek.\Ind{} eat-\textsc{trs}-\Pf{}-\textsc{nmlz} \NegCop{}//
\glft \trsl{The apple is not what Janek ate.} \lingcontext{identificational negation: event presupposed; identification denied}//
\endgl
\xe

\pex
\a
\begingl
\gla Janek lubcí piaščikou staly.//
\glb Janek.\Dir{} apple.\Ind{} eat-\Antip{}-\Pf{}-\textsc{nmlz} \NegCop{}//
\glft \trsl{Janek is not the one who ate the apple.} \lingcontext{agent-identification denied: it was someone else}//
\endgl
\xe

\noindent Example (a) of the first set is compatible with the eating not having
occurred at all. Examples using \ird{staly} presuppose that the eating did occur
and assert that the identified entity is the wrong one: \ird{piaštnikou staly}
says that the apple is misidentified as Janek's eating; \ird{piaščikou staly}
says that Janek is misidentified as the eater. In both cases the implication
that the correct referent exists --- that something else was eaten, or that
someone else ate the apple --- arises from the exhaustivity of the identificational
predication itself rather than from any special property of \ird{staly}.

\ird{Staly} is the same Neg\textsuperscript{0} head here as in plain equative
negation. The contrastive/focus force is a semantic consequence of negating an
identificational (exhaustive) predication, not a distinct grammatical function
of \ird{staly}. See §\,\ref{sec:clausal-negation} for the unified structural
account and §\,\ref{sec:vm-neg-exponents} (\texttt{verb-minim.tex}) for the
formal analysis.

\subsection{The marked copular verb \ird{neh-}: becoming and inchoative
predication}\label{sec:neh}

The root \ird{neh-} provides the overt copular verb. Its paradigm is strongly
suppletive and morphologically irregular; principal forms are given in
Table~\ref{tab:neh-paradigm}.

\begin{table}[ht]
	\footnotesize\sffamily
	\caption{Principal forms of \ird{neh-} (marked copula).}\label{tab:neh-paradigm}
	\medskip
	\begin{tabularx}{0.72\textwidth}{L l}
		\toprule
		{\sc form} & {\sc iridian} \\
		\midrule
		Continuous          & \ird{neha} \\
		Progressive         & \ird{nehime} \\
		Contemplative       & \ird{nehach} \\
		Perfective          & \ird{nek} \\
		Retrospective       & \ird{neaní} \\
		Sequential converb  & \ird{nesu} (colloquial \ird{nu}) \\
		Simultaneous converb & \ird{nešěc} \\
		\bottomrule
	\end{tabularx}
\end{table}

Since neutral equative predication is expressed by zero copula (§\,\ref{sec:zero-copula}), overt \ird{neh-} is semantically marked. It is not used for simple statements of identity or class membership. Instead it signals \term{becoming}, \term{entering a state}, or \term{transition}: the subject is understood as moving from one classification or condition into another.

\pex
\a
\begingl
\gla Marek doktor neha.//
\glb Marek.\Dir{} doctor.\Dir{} \Cop{}-\Cont{}//
\glft \trsl{Marek is becoming a doctor.}//
\endgl
\a
\begingl
\gla Ša byl kapitan nek.//
\glb \Dem{}.\Prox{} child.\Dir{} captain.\Dir{} \Cop{}-\Pf{}//
\glft \trsl{This child became captain.}//
\endgl
\xe

In the continuous aspect, \ird{neh-} may describe a newly entered or recently changed state, with the inchoative flavour modulated by context; a purely stable state assertable without implication of transition calls for zero copula instead. The predicative complement of \ird{neh-} stands in the direct case, as with zero-copula equatives.

The sequential converb \ird{nesu} --- colloquially \ird{nu} --- is a suppletive form established independently of the regular aspect paradigm. Its grammaticalization as a discourse particle is described in Section~\ref{sec:seq-copula} of Chapter~\ref{ch:complex}.

\section{Existential constructions}
\label{sec:existential-constructions}

\subsection{Overview}
\label{sec:existential-overview}%
\label{sec:jec}% retained for cross-reference compatibility
\label{sec:zad}% retained for cross-reference compatibility
\label{sec:jec-zad}% retained for cross-reference compatibility

Iridian expresses existential predication through a pair of suppletive verbal
roots that together constitute the affirmative/negative alternation of a single
functional head, Exist\textsuperscript{0}. The affirmative root \ird{jec-}
asserts the existence or presence of a referent; the negative root \ird{zad-}
denies it. Both inflect for the full aspect system; both are morphologically
irregular and suppletive, with the perfective and retrospective built on
distinct suppletive stems. The two paradigms are given together in
Table~\ref{tab:jec-zad-paradigm}.

\begin{table}[ht]
  \footnotesize\sffamily
  \caption{Principal forms of \ird{jec-} (\textsc{exst}) and \ird{zad-} (\textsc{neg.exst}).}\label{tab:jec-zad-paradigm}
  \medskip
  \begin{tabularx}{0.88\textwidth}{L l l}
    \toprule
    {\sc form} & \ird{jec-} (\textsc{exst}) & \ird{zad-} (\textsc{neg.exst}) \\
    \midrule
    Continuous          & \ird{jeca}   & \ird{zada} \\
    Progressive         & \ird{jecime} & \ird{zdíme} \\
    Contemplative       & \ird{jecach} & \ird{zadach} \\
    Perfective          & \ird{jak}    & \ird{zděk} \\
    Retrospective       & \ird{ján}   & \ird{znení} \\
    \bottomrule
  \end{tabularx}
\end{table}

The two roots are not independent lexical items but are best understood as
suppletive allomorphs of the same Exist\textsuperscript{0} head: \ird{jec-} is
the default exponent, inserted when NegP is absent from the derivation;
\ird{zad-} replaces it when NegP is present, absorbing the negative feature
entirely and leaving Neg\textsuperscript{0} with no independent PF reflex. The
structural and distributional consequences of this analysis are developed in
§\,\ref{sec:neg-exst-alternation} and §\,\ref{sec:vm-jec-zad} below. The
division of labour between \ird{jec-} and the zero-copula construction is
loosely parallel to the Japanese distinction between the existential verbs
\textit{iru} and \textit{aru} on the one hand and the equative
\textit{da}/\textit{desu} on the other: \ird{jec-} covers existence and
presence; zero copula covers identity and class membership.

Iridian has three main productive uses of the existential construction: to
assert existence or presence, including location in a particular place
(\S\,\ref{sec:existence-and-presence}), to express possession
(\S\,\ref{sec:existence-possession}), and certain idiomatic presentational
constructions (\S\,\ref{sec:existential-presentatives}).


\subsection{Existence and presence}\label{sec:existence-and-presence}

The primary use of \ird{jec-} is to assert that a referent exists or is present at the moment of utterance or within a contextually given frame. The entity in question stands in the direct case. \ird{Zad-} forms the direct negative counterpart, asserting non-existence or absence under the same structural conditions.

\pex
\a
\begingl
\gla Byl jeca.//
\glb child.\Dir{} \Exst{}-\Cont{}//
\glft \trsl{There is a child. / A child is present.}//
\endgl
\a
\begingl
\gla Byl zada.//
\glb child.\Dir{} \NegExst{}-\Cont{}//
\glft \trsl{There is no child. / No child is present.}//
\endgl
\xe

\pex
\a
\begingl
\gla Kroumaště na pivo s jeca.//
\glb refrigerator.\Ind{} in beer.\Dir{} still \Exst{}-\Cont{}//
\glft \trsl{There is still (some) beer left in the refrigerator.}//
\endgl
\a
\begingl
\gla Kroumaště na pivo zada.//
\glb refrigerator.\Ind{} in beer.\Dir{} \NegExst{}-\Cont{}//
\glft \trsl{There is no beer in the refrigerator.}//
\endgl
\xe

Within the Minimalist framework, \ird{jec-} is the phonological exponent of a dedicated functional head, \term{Exist}\textsuperscript{0}, which projects \term{ExistP} as the innermost stratum of the clause spine --- the slot occupied by root VP in ordinary verbal clauses. The sole argument it licenses is the \term{existential pivot}: the DP asserted to exist, bearing the direct case as the complement of Exist\textsuperscript{0}. ExistP is then selected by Asp\textsuperscript{0}, which contributes the aspect suffix, and feeds the higher functional sequence including NegP. When NegP is absent, Exist\textsuperscript{0} is spelled out as \ird{jec-}; when NegP is present, suppletive Vocabulary Insertion replaces \ird{jec-} with \ird{zad-}, leaving Neg\textsuperscript{0} with no independent PF exponent. The affirmative and negative structures of \ird{Byl jeca} and \ird{Byl zada} are shown below:

\begin{center}
\begin{forest}
  syntax tree
  [\textsc{AspP}
    [\textit{byl}\\child.\Dir{}\\{\small(pivot)}]
    [\textsc{Asp$'$}
      [\textsc{Asp}\\\textsc{[cont]}: \ird{-a}]
      [\textsc{ExistP}
        [\textsc{Exist$'$}
          [\textsc{Exist}\\\ird{jec-}]
        ]
      ]
    ]
  ]
\end{forest}
\hspace{2em}
\begin{forest}
  syntax tree
  [\textsc{NegP}
    [\textit{byl}\\child.\Dir{}\\{\small(pivot)}]
    [\textsc{Neg$'$}
      [\textsc{Neg}\\{\small($\varnothing$: absorbed by suppletive)}]
      [\textsc{AspP}
        [\textsc{Asp$'$}
          [\textsc{Asp}\\\textsc{[cont]}: \ird{-a}]
          [\textsc{ExistP}
            [\textsc{Exist$'$}
              [\textsc{Exist}\\\ird{zad-}\\{\small(suppletive)}]
            ]
          ]
        ]
      ]
    ]
  ]
\end{forest}
\end{center}

\noindent The pivot surfaces to the left as the absolutive topic constituent, base-generated preverbal under head-finality; no movement is required. In the affirmative it occupies Spec,AspP; in the negative it occupies Spec,NegP. The Neg\textsuperscript{0} feature is discharged by conditioning the suppletive readjustment \ird{jec-} $\rightarrow$ \ird{zad-} at Exist\textsuperscript{0}, and has no residual PF exponent of its own.

\paragraph{Location.} \ird{Jec-} is also used to locate an entity: the entity appears in the direct case and the location is expressed by an indirect noun phrase followed by a postposed clitic adposition. \ird{Zad-} negates location, asserting that the referent is not present at the named location.

\pex
\a
\begingl
\gla Marek dumě na jeca.//
\glb Marek.\Dir{} house.\Ind{} in \Exst{}-\Cont{}//
\glft \trsl{Marek is at home.}//
\endgl
\a
\begingl
\gla Marek dumě na zada.//
\glb Marek.\Dir{} house.\Ind{} in \NegExst{}-\Cont{}//
\glft \trsl{Marek is not at home. / Marek is absent.}//
\endgl
\xe

\noindent The locative phrase is merged as a PP adjunct to ExistP, introducing the spatial domain within which the existential relation holds, without altering the valency of Exist\textsuperscript{0}. The pivot remains the sole direct argument; the structures parallel those for bare existentials above, with the PP adjoined at the ExistP periphery in both affirmative and negative derivations.


\subsection{Possession}\label{sec:existence-possession}

Iridian has no separate verb meaning \trsl{have}. What English presents as \trsl{I have X} is in Iridian construed as \trsl{there is X at me}: the possessor stands in the indirect case, the possessed thing in the direct, and the construction is a special case of existential predication. \ird{Zad-} produces the negative possessive, asserting that the possessed item does not exist within the possessor's sphere.

The ordinary arrangement of the clause places the possessor phrase first, the possessed noun next, and the existential verb last:

\begin{center}
[\emph{possessor}.\Ind{} \ \ird{na}] \ \ [\emph{possessee}.\Dir{}] \ \ \ird{jec-} / \ird{zad-}
\end{center}

\noindent The possessed noun is the thing whose existence is asserted or denied relative to the possessor as reference point. In grammatical terms it is the existential pivot --- the direct argument of Exist\textsuperscript{0} --- while the possessor is an oblique locative domain, a PP adjunct to ExistP, within which the existential relation holds.

\pex
\a
\begingl
\gla Bylam na tóm jeca.//
\glb child.\Ind{} in/on book.\Dir{} \Exst{}-\Cont{}//
\glft \trsl{The child has a book.}//
\endgl
\a
\begingl
\gla Bylam na tóm zada.//
\glb child.\Ind{} in/on book.\Dir{} \NegExst{}-\Cont{}//
\glft \trsl{The child does not have a book.}//
\endgl
\xe

\pex
\a
\begingl
\gla Tóm bylam na jeca.//
\glb book.\Dir{} child.\Ind{} in/on \Exst{}-\Cont{}//
\glft \trsl{As for the book, it is with the child.}//
\endgl
\a
\begingl
\gla Tóm bylam na zada.//
\glb book.\Dir{} child.\Ind{} in/on \NegExst{}-\Cont{}//
\glft \trsl{As for the book, the child does not have it.}//
\endgl
\xe

\noindent The second pair illustrates that topic-initial order is fully available with \ird{zad-} on the same terms as with \ird{jec-}: the alternation belongs to the preverbal field and is independent of the affirmative/negative opposition.

Following Freeze (1992) and the localist tradition, possession is treated here as a structural subtype of location: \trsl{the book exists at the child}. The possessor phrase [\textit{possessor}.\Ind{} \ird{na}] and a plain locative phrase [\textit{location}.\Ind{} \ird{na}] are structurally parallel: both are oblique PPs that modify ExistP without altering the valency of Exist\textsuperscript{0}. The affirmative and negative possessive structures are:

\begin{center}
\begin{forest}
  syntax tree
  [\textsc{ExistP}
    [\textit{PP}\\{\ird{bylam na}}\\{\small(possessor)}]
    [\textsc{ExistP}
      [\textit{tóm}\\book.\Dir{}\\{\small(pivot)}]
      [\textsc{Exist$'$}
        [\textsc{Exist}\\\ird{jec-}]
      ]
    ]
  ]
\end{forest}
\hspace{2em}
\begin{forest}
  syntax tree
  [\textsc{NegP}
    [\textit{tóm}\\book.\Dir{}\\{\small(pivot)}]
    [\textsc{Neg$'$}
      [\textsc{Neg}\\{\small($\varnothing$)}]
      [\textsc{AspP}
        [\textsc{Asp$'$}
          [\textsc{Asp}\\\textsc{[cont]}: \ird{-a}]
          [\textsc{ExistP}
            [\textit{PP}\\{\ird{bylam na}}\\{\small(possessor)}]
            [\textsc{ExistP}
              [\textsc{Exist$'$}
                [\textsc{Exist}\\\ird{zad-}\\{\small(suppletive)}]
              ]
            ]
          ]
        ]
      ]
    ]
  ]
\end{forest}
\end{center}

\noindent The structural parallelism between locative and possessive existentials is a direct consequence of this unified analysis: both constructions merge a PP at the ExistP periphery, and the only grammatical difference is the semantic content of the oblique phrase itself. The suppletive readjustment applies identically in both frames; the PP adjunct is unaffected by the presence or absence of NegP.

Since possession is expressed through the existential verb, possessive clauses naturally inflect across the full aspectual system, and \ird{jec-} and \ird{zad-} participate symmetrically in these aspectual contrasts. The continuous is the common unmarked form: \ird{Bylam na tóm jeca} \trsl{the child has a book}; \ird{Bylam na tóm zada} \trsl{the child does not have a book}. The progressive conveys a relation in the course of arising: \ird{Bylam na tóm jecime} \trsl{the child is getting a book}. The contemplative projects into prospect: \ird{Bylam na tóm jecach} \trsl{the child will have a book}. The perfective, applied to so stative a predicate, yields an inchoative sense: \ird{Bylam na tóm jak} \trsl{the child got a book}. The retrospective places the relation in a prior frame: \ird{Bylam na tóm žaní} \trsl{the child had a book}. In all these forms the construction remains fundamentally existential: the possessed noun remains the direct pivot, the possessor remains oblique, and the aspect alters only the speaker's view of the relation.

Since the possessed noun is the existential pivot, indefinite possessees are especially natural. Definite possessees are not excluded, but with them the clause more readily shades toward a meaning of location, availability, or custody rather than settled ownership --- a consequence of the existential source of the construction.


\subsection{The presentational construction}\label{sec:existential-presentatives}

The continuous form \ird{jeca} licenses a \term{presentational} construction in which the existential pivot is a complex DP modified by a nominalized participial attributive. This construction introduces a referent into discourse by anchoring it to an event in which that referent is the salient non-agentive participant, rather than reporting the event itself. Its canonical shape is:

\begin{center}
[\textit{internal-arg.\Ind{}} \ \textit{participant-oriented-attributive} \ \textit{head-NP}.\Dir{}] \ \ \ird{jeca}
\end{center}

\noindent The modifier is canonically oriented toward the participant the embedded stem privileges: the patient or theme of a transitive predicate, the recipient or beneficiary of an applicative stem, or more generally the non-agentive participant that emerges from the event as the discourse-relevant referent. The construction therefore introduces \trsl{a letter that the woman wrote} or \trsl{a child someone fed} rather than asserting the writing or feeding for its own sake.

\ird{Jeca} is fixed in this use. Temporal and aspectual interpretation come from the embedded participant-oriented modifier, which may be perfective, retrospective, continuous, progressive, or contemplative. The matrix \ird{jeca} signals only that the referent is being introduced into the current discourse.

\pex
\a
\begingl
\gla Obrohnikovní senica jeca.//
\glb read-\textsc{trs}-\Pf{}-\textsc{attr} news.\Dir{} \Exst{}-\Cont{}//
\glft \trsl{There is some news that (I) read.} \lingcontext{patient/theme introduced}//
\endgl
\a
\begingl
\gla Že dbornikovní kíl jeca.//
\glb \textsc{1sg}.\Ind{} hear-\textsc{trs}-\Pf{}-\textsc{attr} song.\Dir{} \Exst{}-\Cont{}//
\glft \trsl{There is a song that I heard.} \lingcontext{stimulus introduced via perception event}//
\endgl
\a
\begingl
\gla Vednikovní pres jeca.//
\glb see-\textsc{trs}-\Pf{}-\textsc{attr} man.\Dir{} \Exst{}-\Cont{}//
\glft \trsl{There is a man that (he) saw.} \lingcontext{perceived referent introduced via completed event}//
\endgl
\a
\begingl
\gla Verce blětébovní byl jeca.//
\glb toy.\Ind{} give-\Appl{}-\Pf{}-\textsc{attr} child.\Dir{} \Exst{}-\Cont{}//
\glft \trsl{There is a child that he gave a toy to.} \lingcontext{recipient introduced via transfer event}//
\endgl
\xe

\noindent The structural basis for this construction is that the pivot position of ExistP may be occupied by any DP of appropriate referential properties. When the pivot is a complex DP --- a head noun modified by a participant-oriented attributive --- ExistP introduces into discourse a referent identified and anchored via an event in which that referent is the salient non-agentive participant. The attributive is an adjunct within the pivot DP; it is \emph{not} a predicate of the matrix clause and contributes no independent propositional content at the matrix level. The matrix Exist\textsuperscript{0} asserts only $\textsc{exist}(x)$ for the complex, event-anchored pivot $x$. This decomposition directly accounts for the fixedness of \ird{jeca}: because the existential head's semantic contribution is reduced to discourse introduction, the least-marked aspect (continuous, unbounded) is the appropriate exponent; any other aspect would impose a viewpoint on the existential state itself, conflicting with the temporal interpretation already contributed by the embedded attributive.

\begin{center}
\begin{forest}
  syntax tree
  [\textsc{AspP}
    [\textsc{DP}\\{\small(pivot)}
      [\textit{AttrP}\\\ird{obrohnikovní}\\{\small read-\textsc{trs}-\Pf{}-\textsc{attr}}]
      [\textit{NP}\\\ird{senica}\\news.\Dir{}]
    ]
    [\textsc{Asp$'$}
      [\textsc{Asp}\\\ird{jec-a}\\{\small(Exist+Cont)}]
      [\textsc{ExistP}
        [\textsc{Exist$'$}
          [\textsc{Exist}\\\ird{jec-}]
        ]
      ]
    ]
  ]
\end{forest}
\end{center}

Because the function of this construction is presentational, agentive heads are disfavoured. An agent is ordinarily presented by a plain verbal clause or, where identification is intended, by an identificational equative (§\,\ref{sec:ident-equative}); the existential-presentational construction is instead the natural means of introducing a participant that has been perceived, affected, selected, or otherwise brought into discourse by the event.

\ex
\begingl
\gla ?Senicět obroždkovní chrás jeca.//
\glb news.\Ind{}.\Def{} read-\textsc{antip}-\Pf{}-\textsc{attr} person.\Dir{} \Exst{}-\Cont{}//
\glft \trsl{There is a person who read the news.} \lingcontext{formally possible, but marked: agentive head disfavoured}//
\endgl
\xe

\noindent This disfavour is structurally grounded: existential predicates cross-linguistically require their pivot to be discourse-new and referentially non-presupposed (the Iridian analogue of Milsark's 1974 definiteness restriction). Agents, as the prototypical topical and volitional participants, are characteristically presupposed or previously introduced, making them poor existential pivots. Patients, themes, recipients, and other non-agentive participants are characteristically discourse-new and are therefore the natural heads of a presentational pivot DP.

The negative counterpart of the presentational construction is formed with \ird{zada}. The negative presentational denies the existence of the event-anchored discourse referent rather than negating the event itself: NegP scopes over ExistP, and the negative asserts that no referent satisfying the event-anchored description exists. As with \ird{jeca}, the continuous form \ird{zada} is fixed in this use; the aspectual interpretation of the embedded event comes entirely from the attributive.

\pex
\a
\begingl
\gla Obrohnikovní senica zada.//
\glb read-\textsc{trs}-\Pf{}-\textsc{attr} news.\Dir{} \NegExst{}-\Cont{}//
\glft \trsl{There is no news that (I) read.} \lingcontext{negative presentational; event-anchored referent denied}//
\endgl
\a
\begingl
\gla Verce blětébovní byl zada.//
\glb toy.\Ind{} give-\Appl{}-\Pf{}-\textsc{attr} child.\Dir{} \NegExst{}-\Cont{}//
\glft \trsl{There is no child that he gave a toy to.} \lingcontext{negative presentational; recipient referent denied}//
\endgl
\xe

\begin{center}
\begin{forest}
  syntax tree
  [\textsc{NegP}
    [\textsc{DP}\\{\small(pivot)}
      [\textit{AttrP}\\\ird{obrohnikovní}\\{\small read-\textsc{trs}-\Pf{}-\textsc{attr}}]
      [\textit{NP}\\\ird{senica}\\news.\Dir{}]
    ]
    [\textsc{Neg$'$}
      [\textsc{Neg}\\{\small($\varnothing$)}]
      [\textsc{AspP}
        [\textsc{Asp$'$}
          [\textsc{Asp}\\\ird{zad-a}\\{\small(NegExst+Cont)}]
          [\textsc{ExistP}
            [\textsc{Exist$'$}
              [\textsc{Exist}\\\ird{zad-}\\{\small(suppletive)}]
            ]
          ]
        ]
      ]
    ]
  ]
\end{forest}
\end{center}

\noindent The Exist\textsuperscript{0} head now asserts $\neg\textsc{exist}(x)$ for the complex pivot: no referent satisfying the event-anchored description exists in the discourse. The attributive remains an adjunct within the pivot DP with no matrix-level propositional force. The disfavour of agentive pivot heads carries over to \ird{zada} by the same definiteness-restriction logic.

The plural particle \ird{ne} cannot be combined with a \ird{zad-} clause; see Section~\ref{sec:definiteness} of Chapter~\ref{ch:nouns}.


\subsection{Suppletive exclusivity: \ird{jec-}, \ird{zad-}, and the NegP alternation}\label{sec:neg-exst-alternation}%
\label{sec:vm-jec-zad}%
\label{sec:vm-jec}

\ird{Jec-} and \ird{zad-} are in strict complementary distribution conditioned by the presence or absence of NegP in the derivation. Wherever NegP is absent, Exist\textsuperscript{0} is spelled out as \ird{jec-}; wherever NegP is present, it is spelled out as \ird{zad-} and \ird{jec-} is excluded. This is an instance of the same conditioned allomorphy that governs Neg\textsuperscript{0}'s exponence across all Iridian clause types: \ird{zá-} (prefix) in verbal clauses, \ird{staly} (free morpheme) in equative clauses, and suppletive root replacement in existential clauses. The three are not separate lexical negation strategies but conditioned phonological reflexes of a single NegP head.

The semantic content of the alternation is straightforwardly compositional. \ird{Jec-} contributes $\textsc{exist}(x)$; \ird{zad-} contributes $\neg\textsc{exist}(x)$, where the negation is the propositional negation licensed by NegP. The same opposition is preserved across all construction types --- bare existential, locative, possessive, and presentational --- and there is no grammatical asymmetry between the two roots in terms of which construction types each permits.

The sole constructional asymmetry is pragmatic. In the presentational use, both \ird{jeca} and \ird{zada} are fixed continuous forms, but the positive presentational functions as a discourse-introducing device (it asserts and presents a new referent), whereas the negative presentational presupposes the availability of the event-anchored description in order to deny a referent matching it. This follows from general properties of existential negation and requires no grammatical stipulation.

\paragraph{The impossibility of \ird{*zájec-}.} A natural question is whether the NEG prefix \ird{zá-} can be used with \ird{jec-} to form a negative existential. It cannot. Forms like \ird{*zájeca} and \ird{*zájecime} are ungrammatical across all existential construction types; negative existential predication is always and exclusively expressed by \ird{zad-}.

\pex
\a
\begingl
\gla Bylam na tóm zada.//
\glb child.\Ind{} in/on book.\Dir{} \NegExst{}-\Cont{}//
\glft \trsl{The child does not have a book.} \lingcontext{\ird{zad-} obligatory in the negative possessive}//
\endgl
\a
\begingl
\gla {*}Bylam na tóm zájeca.//
\glb child.\Ind{} in/on book.\Dir{} \Neg{}-\Exst{}-\Cont{}//
\glft \trsl{(intended: the child does not have a book)} \lingcontext{\ird{zá-} cannot prefix \ird{jec-}}//
\endgl
\xe

\pex
\a
\begingl
\gla Obrohnikovní senica zada.//
\glb read-\textsc{trs}-\Pf{}-\textsc{attr} news.\Dir{} \NegExst{}-\Cont{}//
\glft \trsl{There is no news that (I) read.} \lingcontext{\ird{zad-} obligatory in the negative presentational}//
\endgl
\a
\begingl
\gla {*}Obrohnikovní senica zájeca.//
\glb read-\textsc{trs}-\Pf{}-\textsc{attr} news.\Dir{} \Neg{}-\Exst{}-\Cont{}//
\glft \trsl{(intended: there is no news that (I) read)} \lingcontext{\ird{zá-} cannot prefix \ird{jec-}}//
\endgl
\xe

The structural basis is straightforward. When Neg\textsuperscript{0} is present in an existential clause, it triggers suppletive Vocabulary Insertion at the Exist\textsuperscript{0} node, replacing \ird{jec-} with \ird{zad-}. The Neg\textsuperscript{0} feature is thereby discharged structurally in full; no NegP head remains to be spelled out as \ird{zá-}. The prefix \ird{zá-} is the PF reflex of Neg\textsuperscript{0} only when Neg\textsuperscript{0} has not already been discharged through suppletive readjustment. In verbal clauses, where no suppletive alternation is available at the predicate root, \ird{zá-} is the sole exponent; in existential clauses, the suppletive root exhausts the Neg\textsuperscript{0} feature and \ird{zá-} accordingly never surfaces.

The form \ird{*zájeca} would require Neg\textsuperscript{0} to be doubly exponed: once as the suppletive trigger replacing \ird{jec-} at Exist\textsuperscript{0}, and again as a prefix attaching to the surviving positive root. Since a single syntactic head cannot discharge twice, this configuration is structurally ill-formed. The same logic excludes \ird{*zázada}: stacking \ird{zá-} onto the already-negative \ird{zad-} would require two independent Neg\textsuperscript{0} heads in the same clause, violating the uniqueness of the NegP projection. The prohibition on \ird{zá-} in existential clauses and the prohibition on \ird{*zázad-} are therefore the same prohibition, both following from the single-NegP constraint on the clause spine.
\section{Questions}\label{sec:questions}

\subsection{Yes-no questions}

% [To be written: polar question formation — intonation, question particle]

\subsection{Content questions}

Content questions (\textit{wh}-questions) require the interrogative to occupy the topic (initial) position of the clause. \textit{Wh}-fronting is obligatory in Iridian; in-situ interrogatives are not grammatical in standard Iridian.

% [Port from 2020 archive §\,sec:content-questions; update case labels]

\subsection{Tag, echo, and indirect questions}

% [To be written]

\section{Negation in clausal context}\label{sec:clausal-negation}

At the clausal level, negation in Iridian is expressed by a single functional
projection, NegP, whose Neg\textsuperscript{0} head occupies position 4 in the
clause spine (above AspP, below IP/MoodP; see Table~\ref{tab:th-clause-spine}
in §\,\ref{sec:th-clause-structure}). The surface form of Neg\textsuperscript{0}
varies by predicate type, yielding three distinct phonological reflexes that
together cover all clause types.

\subsection{Verbal clauses: \ird{zá-}}

In finite verbal clauses, Neg\textsuperscript{0} is spelled out as the prefix
\ird{zá-} (reduced to \ird{za-} in unstressed environments). Because the
verbal complex is clause-final under head-finality, and because \ird{zá-}
has no independent prosodic weight, it procliticizes to the left edge of the
verbal complex. The NegP head is syntactically above AspP, but PF
cliticization places the prefix before the root --- a phonological rather
than a syntactic inversion. See §\,\ref{sec:negation} of Chapter~\ref{ch:verbs}
and §\,\ref{sec:vm-neg} (\texttt{verb-minim.tex}) for the full analysis.

\subsection{Equative clauses: \ird{staly}}

In zero-copula equative and class-membership clauses, no verbal complex is
present, and \ird{zá-} therefore has no host to cliticize onto.
Neg\textsuperscript{0} is instead spelled out as the free morpheme \ird{staly},
which surfaces in the canonical head-final position of a NegP head: clause-finally,
after the topic DP and the predicative complement DP. This is the position the
clause spine directly predicts. \ird{Staly} is morphologically invariant; it
bears no aspect, mood, or person marking of any kind, consistent with its status
as a functional head rather than a verbal root. See §\,\ref{sec:zero-copula} for
exemplification of predicational equative negation (\ird{Marek doktor staly})
and §\,\ref{sec:ident-equative} for identificational equative negation
(\ird{Lobek Jancí piaštnikou staly}).

\subsection{Identificational negation: scope contrast with \ird{zá-}}

The availability of participant nominal predicates (§\,\ref{sec:ident-equative})
introduces a structurally motivated contrast between two types of clausal negation
that share the same morphological means but differ in semantic scope:

\begin{description}[leftmargin=2em, labelwidth=!, labelindent=0em]
\item[\upshape Propositional negation (\ird{zá-}).] The prefix \ird{zá-}
  negates the event description expressed by the finite verb. The negated
  clause denies that the event occurred (or was performed by the relevant
  participant). No presupposition is generated about whether the event
  occurred in some other form.
\item[\upshape Identificational negation (\ird{staly} + participant nominal).]
  The combination of a participant nominal predicate and \ird{staly} negates
  the identification: the topic DP is denied as the unique satisfier of the
  nominalized description. The event is presupposed to have occurred; only
  the identity of the participant is in question.
\end{description}

\pex
\a
\begingl
\gla Lobek Janku zápiaštnik.//
\glb apple.\Dir{} Janek.\Trs{} \Neg{}-eat-\textsc{trs}-\Pf{}//
\glft \trsl{Janek did not eat the apple.}
      \lingcontext{propositional: the eating may not have occurred at all}//
\endgl
\a
\begingl
\gla Lobek Jancí piaštnikou staly.//
\glb apple.\Dir{} Janek.\Ind{} eat-\textsc{trs}-\Pf{}-\textsc{nmlz} \NegCop{}//
\glft \trsl{The apple is not what Janek ate.}
      \lingcontext{identificational: eating presupposed; the apple is the wrong identification}//
\endgl
\xe

\noindent The contrastive implication (``it was something else'') in (b) arises
because identificational predications are exhaustive: asserting that $x$ is the
entity satisfying the description entails that no other entity does. Negating
an exhaustive assertion therefore entails the existence of a different satisfier.
This implication is pragmatically cancellable in principle but is the strongly
expected interpretation in discourse.

\subsection{Existential clauses: \ird{zad-}}

In existential clauses, the presence of Neg\textsuperscript{0} in the derivation
triggers suppletive root replacement: the existential root \ird{jec-} is replaced
by \ird{zad-}, which carries the full aspect paradigm of a finite verb.
Neg\textsuperscript{0} itself has no independent PF exponent in this
configuration; it is discharged structurally by licensing the suppletive Vocabulary
Item. The prefix \ird{zá-} does not appear because there is no
Neg\textsuperscript{0} head remaining to be spelled out separately. See
§\,\ref{sec:zad} for paradigm and exemplification.

\subsection{Mood-internal negation: prohibitive and negative hortative}

The prohibitive suffix \ird{-éma} and the negative hortative suffix \ird{-ku}
encode [+\textsc{neg}] as a feature of the MoodP head itself. When MoodP bears
inherent negative force, NegP is not projected, and \ird{zá-} accordingly does
not appear. The negative force of these forms is therefore not a consequence of
NegP but of a higher projection. See §\,\ref{sec:mood} of Chapter~\ref{ch:verbs}.

\subsection{Summary}

The distribution is summarized in the following table:

\begin{table}[ht]
  \footnotesize\sffamily
  \caption{NegP exponents by predicate type.}\label{tab:neg-exponents}
  \medskip
  \begin{tabularx}{\textwidth}{L L L L}
    \toprule
    \textsc{predicate type} & \textsc{affirmative} & \textsc{negative (Neg\textsuperscript{0} exponent)} & \textsc{negation type} \\
    \midrule
    Verbal (indicative)         & verbal complex                & \ird{zá-} prefix on verbal complex        & propositional \\
    Equative / zero copula      & $\emptyset$ copula            & \ird{staly} (free morpheme, clause-final) & predicational \\
    Identificational (part. nom.) & part. nominal + $\emptyset$ & part. nominal + \ird{staly}               & identificational \\
    Existential                 & \ird{jec-}                    & \ird{zad-} (suppletive root)              & propositional \\
    Command (prohibitive)       & \ird{-ím} (imperative)        & \ird{-éma} (NegP not projected)           & --- \\
    Suggestion (hortative)      & \ird{-ká} (hortative)         & \ird{-ku} (NegP not projected)            & --- \\
    \bottomrule
  \end{tabularx}
\end{table}

\noindent The formal analysis, including the DM account of the prefix paradox
(\ird{zá-} as a procliticized NegP head in head-final context) and the full
argument for a unified NegP across all three predicate types, is in
§\,\ref{sec:vm-neg} and §\,\ref{sec:vm-neg-exponents} of the verbal morphosyntax
companion (\texttt{util/stub/verb-minim.tex}).

\section{Comparative and superlative constructions}\label{sec:comparison}

Comparison in Iridian is expressed analytically using particles that combine
with stative verbs. The subject of comparison stands in the direct case; the
standard of comparison takes the indirect case. Two particles are
distinguished: \ird{stá} for direct comparison (\textit{more\ldots than},
\textit{-er than}) and \ird{zrc} for equative comparison (\textit{as\ldots
as}).

\subsection{Direct comparison}

Direct comparison is formed with the particle \ird{stá} placed between the
standard of comparison and the stative verb. The stative verb inflects for the
full range of aspects, giving the usual range of meanings --- ongoing state,
entry into the state, past state, and so on --- within the comparative frame:

\pex
\a
\begingl
\gla Janek Marcí stá četna.//
\glb Janek.\Dir{} Marek.\Ind{} \textsc{comp} be.tall-\Cont{}//
\glft \trsl{Janek is taller than Marek.}//
\endgl
\a
\begingl
\gla Janek Marcí stá četnime.//
\glb Janek.\Dir{} Marek.\Ind{} \textsc{comp} be.tall-\Prog{}//
\glft \trsl{Janek is becoming taller than Marek.}//
\endgl
\a
\begingl
\gla Janek Marcí stá četnik.//
\glb Janek.\Dir{} Marek.\Ind{} \textsc{comp} be.tall-\Pf{}//
\glft \trsl{Janek became taller than Marek.}//
\endgl
\xe

\noindent where \ird{četn-} is \trsl{be tall}. The particle \ird{stá} may also
combine with the \ird{-á} attributive form (§\ref{sec:stative-attr}) to form
an attributive comparative. The particle precedes the \ird{-á} form
prenominally:

\pex
\a \ird{stá četná pres} \trsl{the taller man}
\a \ird{Marcí stá četná pres} \trsl{the man who is taller than Marek}
\xe

\subsection{Equative comparison}

Equative comparison is formed with the particle \ird{zrc}. The structure is
otherwise identical to the direct comparative, with the subject in the direct
case and the standard in the indirect. Negation of the equative falls on the
stative verb itself, which takes the regular negative prefix \ird{zá-}
(§\ref{sec:negation}):

\pex
\a
\begingl
\gla Janek Marcí zrc četna.//
\glb Janek.\Dir{} Marek.\Ind{} \textsc{eq} be.tall-\Cont{}//
\glft \trsl{Janek is as tall as Marek.}//
\endgl
\a
\begingl
\gla Janek Marcí zrc záčetna.//
\glb Janek.\Dir{} Marek.\Ind{} \textsc{eq} \Neg{}-be.tall-\Cont{}//
\glft \trsl{Janek is not as tall as Marek.}//
\endgl
\xe
